% ============================================================
%  Hp3 manuscript -- Supplementary Information
% ------------------------------------------------------------
%  Independent SI document (compiles standalone, like cover_letter.tex).
%  Figure numbering is prefixed "S" (S1, S2, ...). Figures referenced by
%  relative path into the gitignored Dropbox figure tree, as in body.tex.
%  Build:  latexmk supplementary.tex      (xelatex via .latexmkrc engine)
% ============================================================
\documentclass[11pt,a4paper]{article}

% ---- engine / fonts (xelatex; Times New Roman, matching the main text) ----
\usepackage{iftex}
\ifPDFTeX
  \usepackage[T1]{fontenc}\usepackage[utf8]{inputenc}\usepackage{textcomp}
\else
  \usepackage{fontspec}
  \defaultfontfeatures{Scale=MatchLowercase}
  \setmainfont{Times New Roman}
\fi

\usepackage[a4paper,margin=1in]{geometry}
\IfFileExists{microtype.sty}{\usepackage{microtype}}{}
\usepackage{setspace}\onehalfspacing
\setlength{\parindent}{0pt}\setlength{\parskip}{6pt plus 2pt minus 1pt}

\usepackage{graphicx}
\usepackage{caption}
\usepackage{float}
\captionsetup{labelfont=bf,font=small}

% ---- supplementary figure/table numbering: S1, S2, ... ----
\renewcommand{\thefigure}{S\arabic{figure}}
\renewcommand{\thetable}{S\arabic{table}}

\title{\bfseries Supplementary Information\\[4pt]
\large A deeply divergent \emph{Aeromonas} bacteriophage rescues
largemouth bass from lethal \emph{Aeromonas hydrophila} infection}
\date{}

\begin{document}
\maketitle

\begin{figure}[H]
\centering
\includegraphics[width=\linewidth,keepaspectratio]{../analyses/data/121-figs/vc3.png}
\caption{\textbf{Protein-sharing network of phage Hp3 and reference
prokaryotic viruses (vConTACT3).} Each node is a viral genome; nodes are
joined when they share at least one protein cluster, and only viruses
connected to Hp3 are shown. Hp3 (green) groups closely with
\emph{Aeromonas} phage HJ05 (linked by a red edge) and, more distantly,
with \emph{Rhizobium} phage vB\_RleS\_L338C, together forming the deeply
divergent Hp3--HJ05 genus. This cluster is clearly separated from
established families, including its nearest neighbours in the
\emph{Peduoviridae} and \emph{Straboviridae}. Node colour denotes the
family-level assignment given in the legend; ``Other unclassified''
(grey) marks genomes that currently lack a family-level classification.
The network was computed with vConTACT3 against the NCBI viral RefSeq
database (v230) and visualized in Cytoscape.}
\label{fig-vc3-network}
\end{figure}

\end{document}
